% !TEX root = ../main.tex
\section{Problema 1: Sudoku}

\subsection{Descripción del Problema y Abstracción CSP}
% Breve descripción del problema del Sudoku de 9x9 con subcuadrículas de 3x3.
% Justificación de la adecuación del modelo CSP al problema planteado.

\subsubsection{Parámetros y Datos de Entrada}
% Definir matrices o arreglos con las pistas iniciales (números prefijados).

\subsubsection{Variables de Decisión y Dominios}
% Definir con precisión las variables del modelo (ej. matriz X de 9x9).
% Indicar el dominio de cada variable (ej. D(X_ij) = {1, ..., 9}).

\subsubsection{Restricciones del Modelo}
% Formalizar en lenguaje matemático/lógico claro:
% 1. Restricción de filas (todos diferentes).
% 2. Restricción de columnas (todos diferentes).
% 3. Restricción de subcuadrículas 3x3 (todos diferentes).
% 4. Asignación de pistas iniciales.

\subsection{Detalles de Implementación (\texttt{sudoku.mzn})}
% Describir aspectos relevantes de la implementación sin incluir todo el código.

\subsubsection{Diferentes Implementaciones de Restricciones}
% Comparar implementaciones: ej. alldifferent vs pares de desigualdades (!=).

\subsubsection{Restricciones Redundantes y Rompimiento de Simetrías}
% Explicar si aplican restricciones redundantes o si hay simetrías. En caso negativo, argumentar por qué.

\subsection{Estrategias de Búsqueda y Distribución}
% Describir al menos dos anotaciones de búsqueda exploradas en MiniZinc.
% Ejemplo: int_search(..., first_fail, indomain_min, complete), dom_w_deg, etc.

\subsection{Pruebas y Análisis Experimental}
% Presentar la batería de pruebas (ej. instancias fáciles, medias, difíciles).
% Incluir tablas o gráficas con estadísticas de MiniZinc (nodos explorados, nodos fallidos, tiempo de respuesta).

\begin{table}[htbp]
    \centering
    \begin{tabular}{lcccc}
        \toprule
        Instancia & Estrategia & Tiempo (s) & Nodos explorados & Nodos fallidos \\
        \midrule
        Instancia 1 & Default & -- & -- & -- \\
        Instancia 1 & Custom 1 & -- & -- & -- \\
        Instancia 2 & Default & -- & -- & -- \\
        \bottomrule
    \end{tabular}
    \caption{Comparación de rendimiento para el problema del Sudoku.}
    \label{tab:sudoku_resultados}
\end{table}

\subsection{Conclusiones del Ejercicio}
% Análisis crítico de los resultados, ventajas y desventajas de las implementaciones ensayadas.
